\section{Conclusion}\label{sec:conclusion}

Public climate protection and private production resilience are jointly determined. A safer jurisdiction can attract primary production, but the same protection can reduce firms' willingness to maintain costly geographic backup capacity. When local governments compete for mobile production, they value the relocation benefit even if the induced reduction in private redundancy has already made additional public protection socially undesirable.

A simple two-jurisdiction oligopoly model formalizes this interaction. Firms choose primary locations and backup readiness before localized disasters, and surviving firms compete in the product market. The model generates a nonempty open set of primitives in which the coordinated planner chooses no additional local protection while decentralized jurisdictions choose strictly positive protection. The policy conflict is absent from the relevant nested benchmarks: fixed locations remove the attraction motive, while fixed redundancy removes the organizational response that changes the planner's valuation of protection.

The broader implication is methodological. Place-based resilience policy should not be evaluated as an engineering change in local hazard exposure holding private organization fixed. Where firms can move production and redesign contingency capacity, public adaptation changes the industrial resilience architecture itself. Decentralization then adds a second wedge because local governments value where production is located, not only whether the national production system remains resilient.
